\documentclass[11pt]{article}
\usepackage[a4paper,margin=1in]{geometry}
\usepackage{amsmath,amssymb}
\usepackage{enumitem}
\usepackage[colorlinks=true,linkcolor=blue,urlcolor=blue,citecolor=blue]{hyperref}
\usepackage{titlesec}

% Section headers
\newcommand{\secx}[1]{\section{#1}}
\newcommand{\subsecx}[1]{\subsection{#1}}

% Footnote symbols (*, dagger, ...) instead of numbers, standard for author affiliations
\renewcommand{\thefootnote}{\fnsymbol{footnote}}

\title{\textbf{Contract Playbook Comparison System}\\[2pt]
\large A Retrieval-Augmented System for Automated Contract Clause Review}
\author{
Mannat, CSED\thanks{Roll No: 1024240023, Email: mannat\_be24thapar.edu}
\and
Upasna Wadhwa, CSED\thanks{Roll No: 1024240022, Email: uwadhwa\_be24thapar.edu}
\and
Yuvicka, CSED\thanks{Roll No: 1024240016, Email: yuvicka\_be24thapar.edu}
}
\date{August 17, 2026}

\begin{document}

\begin{center}
{\large\bfseries Contract Playbook Comparison System}\\[4pt]
{\normalsize Project Proposal for UCS503}\\[2pt]
{\normalsize Submitted to: Dr.\ Jeelani Asif}\\[2pt]
{\normalsize Thapar Institute of Engineering and Technology}
\end{center}

\vspace{6pt}

\begin{center}
Mannat, DSAI\footnotemark[1] \qquad\qquad
Upasna Wadhwa, DSAI\footnotemark[2] \qquad\qquad
Yuvicka, DSAI\footnotemark[3]
\end{center}

\begin{center}
August 16, 2026
\end{center}

\footnotetext[1]{Roll No: 1024240023, Email: mmannat\_be24@thapar.edu}
\footnotetext[2]{Roll No: 1024240022, Email: uwadhwa\_be24@thapar.edu}
\footnotetext[3]{Roll No: 1024240016, Email: yyuvicka\_be24@thapar.edu}

\vspace{10pt}

\secx{Higher-order goal}
This project focuses on improving the reliability and consistency of corporate legal and
procurement decision-making by minimising the manual effort required to review vendor
contracts against internal risk policy. While the underlying technique (retrieval-augmented
generation over legal text) is broadly applicable, the immediate engineering objective is to translate
grounded, citation-backed clause comparison into a reliable, deployable software pipeline.

\secx{Time-to-value}
\begin{itemize}
  \item We will deliver meaningful updates quickly by prototyping the core
  retrieve-compare-report workflow on a small playbook (5 to 8 clause types) and
  validating it against some contracts within a short iteration window.
  \item We will prioritize fast feedback over perfection by building each pipeline
  stage (ingestion, retrieval, comparison, reporting) as an independently
  testable module, so partial progress may be visible every week.
  \item Success shall be defined by what can be delivered and measured in the first
  iteration that could be a working retrieval step evaluated against ground-truth clause
  spans and then improved based on analysis of retrieval and classification errors.
\end{itemize}

\secx{Problem Statement}
Companies bargain and sign large volumes of vendor, SaaS, and service contracts.
Legal and procurement teams maintain an internal \emph{playbook}: a set of pre-approved
clause positions representing acceptable risk (e.g., minimum liability caps, maximum
termination notice periods, and compulsory auto-renewal opt-outs). When a new contract
arrives, a reviewer must physically read the document (often 20--100+ pages), locate each
relevant clause, compare it against the playbook, and flag variations. Common pain points include:

\begin{itemize}
  \item \textbf{Slow, non-scaling review:} manual review takes much time per contract and
  does not scale with a large number of contracts.
  \item \textbf{Inconsistent outcomes:} different reviewers catch different variations,
  so review quality varies by who reviews it.
  \item \textbf{No audit trail:} verdicts are hardly linked back to the exact source
  text that produced them, making them hard to prove or dispute later.
  \item \textbf{Hallucination risk of naive LLM use:} feeding an entire contract to a
  general-purpose LLM and asking, "Does this match our policy?" produces
  unverifiable, sometimes fabricated, answers  unacceptable when a missed
  liability clause has real financial consequences.
\end{itemize}

\noindent\textbf{Why this matters now:} Contract Lifecycle
Management (CLM) is an established, funded software category (e.g., Ironclad,
LinkSquares, Lawgeex, Evisort), confirming this is a real workflow gap and not a
contrived exercise. A retrieval-grounded approach is well suited here precisely because
the cost of an ungrounded answer is high.

\secx{Proposed Solution}

\subsecx{Overview}
Building a retrieval-augmented pipeline, the \emph{Contract Playbook Comparison System},
with the following parts:
\begin{itemize}
  \item \textbf{Playbook definition:} 5 to 8 clause categories (e.g., Termination,
  Liability Cap, Auto-Renewal, Indemnification, Governing Law) each with a standard documented position.
  \item \textbf{Clause retrieval:} hybrid semantic + keyword search that locates the
  section(s) of an incoming contract mapping it to each playbook clause.
  \item \textbf{Clause comparison:} an LLM classifies each retrieved section as
  \emph{Match}, \emph{Deviation}, or \emph{Missing} according to the playbook
  position, with a natural language explanation.
  \item \textbf{Structured report:} one row per playbook clause, showing status,
  explanation, and an exact citation into the source contract.
\end{itemize}

\subsecx{Core workflow}
\begin{enumerate}
  \item An incoming contract is parsed and chunked into indexable segments.
  \item For each playbook clause type, the retrieval stage will return the top-$k$
  candidate contract chunks.
  \item If no chunk passes a similarity threshold, the clause is classified
  \emph{Missing} directly and hence the LLM is never asked to guess from nothing.
  \item Otherwise, the LLM compares the retrieved text against the playbook position
  and returns \emph{Match} / \emph{Deviation} with a one-sentence explanation.
  \item Results are put into a structured report with source citations.
\end{enumerate}

\subsecx{Operational constraints}
\begin{itemize}
  \item To use a well-established public dataset (CUAD) rather than sourcing and
  labelling proprietary contracts.
  \item Avoid dependence on unproven algorithms rather: favour standard,
  well-documented retrieval and evaluation techniques.
  \item Designing for a reproducible local/demo deployment (Streamlit or Gradio) rather
  than production-grade infrastructure.
\end{itemize}

\secx{Solution Approach}
This is an engineering course project, so we focus on building and validating a
modular pipeline rather than on unusual algorithm research.

\subsecx{Dataset and grounding}
We will use CUAD (Contract Understanding Atticus Dataset): 510 real commercial contracts
with 13,000+ expert-annotated clause spans across 41 categories.
\begin{itemize}
  \item Select 5--8 clause categories with reasonable positive-example counts in
  CUAD and define a documented standard position for each.
  \item Use held-out CUAD contracts as test inputs, giving automatic ground truth for
  retrieval evaluation (does the system find the correct labelled span?).
  \item Hand-label a small set (\textasciitilde50--100) of clause-comparison pairs as
  Match/Variation/Missing for categorising evaluation.
\end{itemize}

\subsecx{Pipeline architecture}
A staged retrieval-augmented pipeline:
\begin{itemize}
  \item \textbf{Ingestion:} PDF/text parsing (pdfplumber/PyMuPDF) and intersecting
  chunking (\textasciitilde300--500 tokens, \textasciitilde15\% overlap) tuned
  for long legal clauses.
  \item \textbf{Indexing:} sentence-transformer embeddings stored in FAISS/ChromaDB,
  alongside a BM25 keyword index for hybrid retrieval,  important because
  exact legal terminology (`"indemnify", "force majeure") often matters
  more than pure semantic similarity.
  \item \textbf{Comparison:} a structured LLM prompt that receives only the playbook
  position and the retrieved chunk(s), prohibiting the model from answering off
  of unretrieved, hallucinated context.
  \item \textbf{Reporting:} a single table with  clause type, status, explanation, and
  citation  as the reviewer-facing output.
\end{itemize}

\subsecx{CI/CD and fast delivery enabler}
CI/CD will be used to:
\begin{itemize}
  \item Automatically run retrieval and classification validation scripts on every
  change to the pipeline.
  \item Produce a repeatable demo build (Streamlit/Gradio) for weekly review.
  \item Enable safe, frequent iteration on every pipeline stage without
  reevaluating the whole system manually.
\end{itemize}

\secx{Evaluation Criterion}
We define success using metrics bind to each pipeline stage, each with its own
ground truth.

\subsecx{Primary evaluation metric}
\textbf{Retrieval accuracy}: Does the system return the correct contract span for each
playbook clause type, measured directly against CUAD's labelled spans?
\begin{itemize}
  \item Target: Recall@$k$ (k=3-5) and mean reciprocal rank (MRR) reported per
  clause category.
  \item Attribution: computed automatically by matching retrieved chunk offsets
  against CUAD ground-truth spans.
\end{itemize}

\subsecx{Secondary evaluation metrics}
\begin{itemize}
  \item \textbf{Classification accuracy:} agreement between the pipeline's
  match/deviation/missing verdicts and the team's manually-labelled test set.
  \item \textbf{Faithfulness:} Whether the LLM's explanation is led by the
  retrieved text (checked manually on a selection and optionally via an
  NLI-based check as in the RAGAS framework).
  \item \textbf{Coverage:} percentage of playbook clause types for which the system
  returns a confident retrieval (above threshold) versus falling back to
  \emph{Missing}.
\end{itemize}

\subsecx{Validation plan}
\begin{itemize}
  \item Evaluating retrieval and classification on a held-out split of CUAD contracts
  not used to design the playbook positions.
  \item Comparing hybrid (semantic + BM25) retrieval against semantic-only retrieval
  as an ablation, to justify the hybrid design choice.
\end{itemize}

\secx{Scalability}
The system should scale in both conceptual and practical senses.
\begin{enumerate}
  \item \textbf{Scalable (conceptually):} the pipeline should support supplementary
  playbook clause types and larger contract volumes by
  \begin{itemize}
    \item decoupling ingestion, retrieval, comparison, and journalising stages,
    \item indexing embeddings for sub-linear nearest-neighbour lookup,
    \item batching LLM comparison calls independently per clause type.
  \end{itemize}
  \item \textbf{Operationally feasible (practically):} the demo should run on
  commodity resources:
  \begin{itemize}
    \item a local or lightly-hosted vector index (FAISS/ChromaDB),
    \item any accessible LLM API or a small open model to control cost,
    \item CI-run evaluation scripts that do not require GPU infrastructure.
  \end{itemize}
\end{enumerate}

\secx{Engine availability heuristic}
A mature set of  "engines" (tooling/libraries) is available to implement this
without inventing a new infrastructure:
\begin{itemize}
  \item Sentence-embedding models (\texttt{sentence-transformers}, e.g.\
  \texttt{all-mpnet-base-v2}, or a legal-domain model such as Legal-BERT).
  \item A vector store (FAISS or ChromaDB) and a keyword index (\texttt{rank\_bm25}).
  \item An accessible LLM API for the comparison/classification step.
  \item An evaluation framework purpose-built for RAG (RAGAS) instead of a
  manual evaluation system.
  \item A lightweight demo UI framework (Streamlit or Gradio).
\end{itemize}
We will use these existing components to focus on pipeline correctness and evaluation
quality rather than infrastructure-building.

\secx{Project scope and deliverables}

\subsecx{Initial deliverable}
\begin{itemize}
  \item Playbook definition for 5--8 clause categories with documented standard
  positions.
  \item Ingestion and chunking pipeline for CUAD contract text.
  \item Embedding index (FAISS/ChromaDB) with basic semantic retrieval.
  \item Retrieval evaluation against CUAD ground-truth spans (Recall@$k$, MRR).
  \item CI: automated evaluation run + linting on every change.
\end{itemize}

\subsecx{Subsequent deliverables}
\begin{itemize}
  \item Hybrid retrieval (BM25 + semantic) and ablation comparison against
  semantic-only retrieval.
  \item LLM-based clause comparison with structured Match/Deviation/Missing output.
  \item Hand-labelled classification test set and accuracy/faithfulness evaluation.
  \item Streamlit/Gradio demo: contract selection $\rightarrow$ playbook comparison
  $\rightarrow$ cited report.
\end{itemize}

\secx{Risks and mitigations}
\begin{itemize}
  \item \textbf{LLM hallucination even with retrieved context:} keep comparison
  prompts strict (``only use the provided text; say `cannot determine' if
  unclear''), and always surface the retrieved source text alongside the
  verdict so it is independently verifiable.
  \item \textbf{Class imbalance in CUAD categories:} check per-category label counts
  in week 1 and choose playbook clause types with enough representation.
  \item \textbf{Fuzzy clause boundaries:} allow retrieval to return multiple adjacent
  chunks rather than a single best match, since clauses can span paragraphs.
  \item \textbf{Domain-specific phrasing confusing embeddings:} mitigate with hybrid
  (semantic + keyword) retrieval, since exact legal terms often carry more
  signal than semantic similarity alone.
\end{itemize}

\secx{Summary}
This project suggests a retrieval-augmented pipeline that automates contract clause
review against an internal playbook, producing grounded, citation-backed verdicts
instead of unverifiable LLM output. It meets the course criteria by decomposing the
system into independently testable modules, defining measurable validation criteria
at each stage (retrieval accuracy, classification accuracy, and faithfulness), and
grounding the whole design in a well-established public legal-NLP benchmark (CUAD)
rather than research-driven novelty.

\end{document}
